\section{Introduction}

Molecular property prediction is a critical component of early-stage drug discovery, enabling the prioritization of candidate compounds before expensive synthesis and biological assays~\cite{vamathevan2019applications}. Machine learning approaches---including graph neural networks~\cite{gilmer2017neural}, fingerprint-based classifiers~\cite{rogers2010extended}, and descriptor-based models---have become standard tools, with MoleculeNet~\cite{wu2018moleculenet} providing widely adopted benchmarks across physicochemical, enzymatic, and toxicity tasks.

More recently, large language models (LLMs) have been applied to molecular reasoning by converting SMILES notation into natural language contexts for text-based inference~\cite{guo2024llm, ross2022large}. Within this line of work, an observation has been reported in prior work: hallucinated molecular descriptions---scientifically styled but factually incorrect text---have been associated with improved downstream prediction performance in specific settings~\cite{hao2024hallucination}. This finding is in tension with the common assumption that strict factual correctness is a prerequisite for useful textual augmentation. We note that zero-shot LLM inference is not intended to replace specialized molecular prediction models; rather, the LLM serves as a controlled experimental substrate---a ``cognitive assay''---for investigating how semantic perturbations propagate through text-conditioned prediction pipelines.

However, existing investigations of this phenomenon share several limitations. First, hallucination is typically treated as a monolithic category, obscuring potentially important differences among types of factual error. Second, prior studies often lack semantic controls capable of disentangling the contribution of molecular-relevant content from generic scientific language effects (stylistic priming). Third, it remains unclear whether these hallucination-associated performance variations generalize across prediction tasks with different underlying knowledge requirements. Understanding this sensitivity is important for interpreting LLM-augmented predictions, particularly in settings where textual augmentation may introduce uncontrolled semantic perturbations.

The primary contribution of this work is the introduction of a controlled semantic perturbation framework to systematically evaluate how structured semantic alterations influence LLM-based molecular prediction. We define a lightweight, heuristic categorization scheme comprising three prompt-constrained perturbation operators---Structural Phantom (SP), Property Inversion (PI), and Mechanism Fabrication (MF)---alongside a residual Contextual Confabulation (CC) class. We note that these labels describe the \textit{prompt design intent}; as detailed in Section~\ref{sec:taxonomy}, the actual outputs may not conform to the intended taxonomy class. We evaluate this framework through a nine-condition ablation on two MoleculeNet benchmarks representing distinct regimes: BBBP (physicochemical) and BACE (enzyme inhibition).

Our key contributions include:
\begin{itemize}
    \item A controlled perturbation framework enabling the systematic isolation of hallucination-associated performance variations through semantic and stylistic disruption.
    \item A lightweight categorization scheme for molecular description errors, operationally defined as perturbation operators to maintain interpretability and reproducibility.
    \item A task-dependent divergence in semantic augmentation sensitivity: prompt-constrained structural augmentation (C4a) yields directional improvements on BBBP, while most hallucination conditions significantly degrade BACE performance. Post-hoc verification reveals 0\% SP classification for C4a, indicating that the prompt's structural topic focus---rather than hallucination type---underlies the observed effect.
    \item Characterization of prediction distribution compression on BACE, where hallucination augmentation shifts bimodal baseline distributions toward non-discriminative narrow peaks.
    \item A random-permutation control experiment establishing that semantic content, rather than stylistic register alone, is the primary contributor to observed performance shifts.
    \item Architecture-dependent hallucination sensitivity, where smaller models exhibit directional performance gains under hallucination augmentation on physicochemical tasks, while larger models show reversed sensitivity.
\end{itemize}

This exploratory study is bounded to two LLM scales and two datasets ($N=204$ for BBBP, $N=152$ for BACE). As a cognitive assay of semantic sensitivity, the findings are intended to structure future empirical investigation into LLM augmentation rather than to establish causal mechanisms.
